\section{Conclusion}
\label{sec:conclusion}

In this paper, we presented a rigorous empirical study of the vulnerabilities of test-time dynamic model merging and routing. Specifically, we analyzed: (1) calibration data scarcity, which leads standard parametric routers to overfit and experience transductive collapse, and (2) stream batch heterogeneity, which leads to batch-averaging representation smoothing and heterogeneity collapse. 

To overcome these failures, we proposed \textbf{Confidence-Gated Hybrid Routing (CGHR)}, a dual-pathway system that gates predictions sample-by-sample, falling back to non-parametric subspace routing (PFSR) when parametric confidence is low. Furthermore, we integrated \textbf{Micro-Batch Homogenization (MBH)}, which dynamically clusters heterogeneous stream batches into isolated, homogeneous sub-batches on the fly.

Through extensive sweeps over 5 independent random seeds, we mapped the complete trade-off curve across gating thresholds, calibration sample sizes, and deployment batch sizes. Our experimental results provide overwhelming empirical proof of our system's advantages:
\begin{itemize}
    \item CGHR secures a robust peak performance envelope at intermediate thresholds, outperforming both pure parametric and non-parametric baselines.
    \item Under severe calibration data scarcity ($N=16$), CGHR achieves stable performance with near-zero seed variance, completely bypassing overfitting.
    \item MBH fully protects dynamic merging against heterogeneity collapse, maintaining flat, robust performance curves up to batch sizes of $B=512$.
\end{itemize}

\subsection{Limitations and Future Work}
\label{sec:limitations}
While our thorough simulation-based audits provide high-signal empirical observations, we explicitly identify three major limitations of the current study that outline essential avenues for future research, which we discuss in exhaustive detail in the Appendices:
\begin{enumerate}
    \item \textbf{Synthetic Sandbox Reliance and Feature Overlap}: Our quantitative evaluations are restricted to a 1-layer synthetic \emph{Isolating Coordinate Sandbox}. In actual pre-trained deep neural networks (e.g., Transformers), representation spaces are highly overlapping and non-orthogonal, causing gating ambiguity and breakdown of the UNC-PFSR Equivalence. To restore the equivalence properties, we propose an elegant mathematical formulation using SVD subspace projection operators to project representations onto task-specific manifolds on-the-fly (see Appendix~\ref{app:svd_projections_detailed}).
    \item \textbf{Adaptive and Dynamic Confidence Gating}: The optimal confidence threshold ($\gamma_{\text{conf}} \approx 0.85$) is searched statically. Under non-stationary, volatile streams, we propose a dynamic self-calibrating threshold $\gamma_{\text{conf}}(t) = \gamma_{\text{base}} + \eta \bar{H}_t$ based on rolling stream entropy to protect the pipeline (see Appendix~\ref{app:limitations_and_future_work}).
    \item \textbf{SVD Projection Scaling to Real-World Benchmarks}: We provide a concrete, step-by-step qualitative roadmap for deploying our SVD subspace projections on real-world multi-task benchmarks (e.g., DomainNet or GLUE) with pre-trained Transformer backbones and LoRA adapters, demonstrating how the singular value decay allows for low-rank, memory-efficient constructs (see Appendix~\ref{app:svd_scaling_roadmap}).
\end{enumerate}